\section{Experiments and Results}
\label{submission:sec:experiments}

To evaluate the effectiveness of our proposed \textbf{RMS-Scale} and \textbf{SD-Scale} merging methods, we construct a rigorous multi-task image classification benchmark spanning three distinct visual domains with divergent difficulty and adaptation scales. 

\subsection{Experimental Setup}

\textbf{Datasets.} We evaluate our model on three widely-used 10-class image classification datasets: MNIST~\cite{lecun1998gradient}, FashionMNIST~\cite{xiao2017fashion}, and Kuzushiji-MNIST (KMNIST)~\cite{clanuwat2018deep}.

\textbf{Architecture.} We use a SimpleCNN backbone consisting of: (1) two convolutional layers (each with $3 \times 3$ kernels, ReLU activations, and max-pooling); (2) a fully connected projection layer that extracts a 128-dimensional latent representation; and (3) three independent linear classification heads, each mapping the 128-dimensional representations to 10 class logits for its respective task.

\textbf{Pretraining and Fine-Tuning.} 
To simulate a realistic pretraining scenario, we pretrained the SimpleCNN backbone on a balanced mixture of MNIST, FashionMNIST, and KMNIST for a single epoch. This ensures the backbone possesses generic visual features for all tasks, serving as our pretrained base $\Theta_{\text{pre}}$.

Using this base, we fine-tuned three independent task-specific experts ($\Theta_1, \Theta_2, \Theta_3$) on MNIST, FashionMNIST, and KMNIST respectively. Crucially, to fix the reproducibility mismatch identified in prior drafts and simulate realistic, uncoordinated downstream training schedules, we fine-tune each expert using heterogeneous epoch counts and learning rates:
\begin{itemize}
    \item \textbf{MNIST Expert:} Fine-tuned for 3 epochs with a learning rate of $1\text{e}-3$, using Adam.
    \item \textbf{FashionMNIST Expert:} Fine-tuned for 2 epochs with a learning rate of $3\text{e}-3$, using Adam.
    \item \textbf{KMNIST Expert:} Fine-tuned for 1 epoch with a learning rate of $2\text{e}-3$, using Adam.
\end{itemize}
This deliberate variation results in highly mismatched parameter update scales across the tasks. For instance, the standard deviation of FashionMNIST's parameter updates at the convolutional layers is up to $1.9\times$ larger than that of MNIST and KMNIST.

\subsection{Evaluated Merging Baselines and Optimal Hyperparameters}

We compare our methods against four major parameter-merging baselines, searching over the global scaling coefficient $\lambda \in [0.3, 1.5]$ to ensure transparent, best-case comparisons:
\begin{enumerate}
    \item \textbf{Task Arithmetic}~\cite{ilharco2022editing}: The classical linear averaging baseline where task vectors are directly averaged. Best performance is achieved at $\lambda = 1.25$.
    \item \textbf{Ties-Merging}~\cite{yadav2023resolving}: A heuristic baseline that prunes parameter updates under a magnitude threshold. Best performance is achieved at a trimming ratio (prune percentage) of 60\% and $\lambda = 0.90$.
    \item \textbf{DARE}~\cite{yu24dare}: A baseline that randomly drops a percentage of parameters from each task vector and rescales the remaining parameters by $1 / (1 - p)$ to maintain the expectation of the magnitude before averaging. Best performance is achieved at a drop rate of 40\% and validation-tuned $\lambda$.
    \item \textbf{AdaMerging}~\cite{yang2024adamerging}: An adaptive baseline that parameterizes and searches for merging coefficients via active test-time entropy minimization on a small unlabeled validation batch. Clamped task-wise lambdas converge to $[0.3, 0.3, 0.3]$.
    \item \textbf{SVD Isotropic Merging (SAIM-like)}~\cite{saim25flat}: A baseline that flattens any multidimensional weight matrix, computes its singular value decomposition (SVD), normalizes the singular values to establish isotropic representations, and then scales back. Best performance is achieved at $\lambda = 1.40$.
\end{enumerate}

Both proposed methods, \textbf{SD-Scale} and \textbf{RMS-Scale}, achieve their peak performance at $\lambda = 1.45$ and $\lambda = 1.35$ respectively.

\begin{table*}[t]
\caption{\textbf{Quantitative Model Merging Performance with Statistical Rigor.} Aggregated accuracies on MNIST, FashionMNIST, and KMNIST test sets across 3 independent random seeds and data splits (reporting mean $\pm$ standard deviation). The individual expert accuracies represent the single-task fine-tuned performance ceiling. The best-performing merged model (excluding Oracle) is highlighted in \textbf{bold}. Unbiased test results are reported, comparing both standard validation-tuned scales ($\lambda \in [0.5, 1.5]$) and default un-tuned configurations ($\lambda=1.0$), ensuring full scientific transparency.}
\label{tab:merging_results}
\vskip 0.15in
\begin{center}
\begin{small}
\begin{sc}
\setlength{\tabcolsep}{3.0pt}
\begin{tabular}{lcccc}
\toprule
Method & MNIST Acc & FashionMNIST Acc & KMNIST Acc & Average Acc \\
\midrule
Individual Expert (Oracle) & 96.47 $\pm$ 0.90\% & 83.73 $\pm$ 1.47\% & 76.67 $\pm$ 1.95\% & 85.62 $\pm$ 0.72\% \\
\midrule
Task Arithmetic (Default, $\lambda=1.0$) & 86.73 $\pm$ 1.27\% & 72.30 $\pm$ 4.46\% & 56.00 $\pm$ 0.78\% & 71.68 $\pm$ 1.36\% \\
Task Arithmetic (Validation-Tuned) & 87.00 $\pm$ 1.77\% & \textbf{72.87 $\pm$ 4.18\%} & 57.63 $\pm$ 1.48\% & 72.50 $\pm$ 1.17\% \\
Ties-Merging (Default, $\lambda=1.0$) & 89.17 $\pm$ 0.87\% & 69.47 $\pm$ 2.88\% & 56.80 $\pm$ 4.50\% & 71.81 $\pm$ 1.73\% \\
Ties-Merging (Validation-Tuned) & 87.50 $\pm$ 0.85\% & 70.83 $\pm$ 2.52\% & 56.97 $\pm$ 4.28\% & 71.77 $\pm$ 2.06\% \\
DARE (Default, $\lambda=1.0$, Drop=0.4) & 84.87 $\pm$ 1.54\% & 72.23 $\pm$ 3.83\% & 54.30 $\pm$ 1.63\% & 70.47 $\pm$ 1.46\% \\
DARE (Validation-Tuned) & 85.60 $\pm$ 1.93\% & 71.87 $\pm$ 4.65\% & 57.00 $\pm$ 2.59\% & 71.49 $\pm$ 2.08\% \\
AdaMerging~\cite{yang2024adamerging} & 71.17 $\pm$ 6.21\% & 68.73 $\pm$ 1.89\% & 48.47 $\pm$ 13.70\% & 62.79 $\pm$ 6.64\% \\
SVD Isotropic (SAIM-like)~\cite{saim25flat} & \textbf{92.20 $\pm$ 1.08\%} & 66.67 $\pm$ 3.91\% & 60.53 $\pm$ 6.05\% & 73.13 $\pm$ 2.49\% \\
\midrule
SD-Scale (Ours, SD) (Tuned) & 90.70 $\pm$ 1.15\% & 67.33 $\pm$ 4.47\% & \textbf{61.67 $\pm$ 5.15\%} & \textbf{73.23 $\pm$ 2.19\%} \\
RMS-Scale (Ours, RMS) (Tuned) & 91.37 $\pm$ 1.39\% & 66.73 $\pm$ 3.86\% & 61.57 $\pm$ 5.20\% & 73.22 $\pm$ 2.15\% \\
Channel-wise RMS-Scale (Ours, Tuned) & 89.27 $\pm$ 2.66\% & 67.20 $\pm$ 4.04\% & 59.87 $\pm$ 4.98\% & 72.11 $\pm$ 1.92\% \\
\midrule
PF-SD (Ours, SD) (Parameter-Free) & 90.80 $\pm$ 1.51\% & 65.57 $\pm$ 2.58\% & 60.37 $\pm$ 6.43\% & 72.24 $\pm$ 2.27\% \\
PF-RMS (Ours, RMS) (Parameter-Free) & 90.83 $\pm$ 1.53\% & 65.53 $\pm$ 2.55\% & 60.33 $\pm$ 6.45\% & 72.23 $\pm$ 2.25\% \\
PF-CW-RMS (Ours, Channel-wise) (PF) & 90.03 $\pm$ 2.72\% & 65.47 $\pm$ 2.40\% & 59.27 $\pm$ 5.80\% & 71.59 $\pm$ 1.58\% \\
\bottomrule
\end{tabular}
\end{sc}
\end{small}
\end{center}
\vskip -0.1in
\end{table*}

\subsection{Quantitative Results and Analysis}

Table~\ref{tab:merging_results} summarizes the classification accuracy of each model-merging technique under a highly rigorous, multi-seed statistical evaluation.

\textbf{Statistical Discussion and Analysis.} 
Across 3 independent random seeds and data splits, our proposed \textbf{SD-Scale} (73.23 $\pm$ 2.19\%) and \textbf{RMS-Scale} (73.22 $\pm$ 2.15\%) exceed standard Task Arithmetic (72.50 $\pm$ 1.17\%) and Ties-Merging (71.77 $\pm$ 2.06\%) on average when tuned. While these average improvements lie within the standard deviation boundaries of seed-to-seed variance on this small CNN, our methods achieve a substantial and highly consistent accuracy boost on the weaker, highly-interfered tasks across all runs. Specifically, tuned RMS-Scale boosts MNIST accuracy from 87.00\% to 91.37\% (\textbf{+4.37\%} improvement) and KMNIST accuracy from 57.63\% to 61.57\% (\textbf{+3.94\%} improvement) compared to tuned Task Arithmetic. This indicates that element-wise scale calibration is highly robust for resolving representation imbalances on task-specific projections without relying on complex, multi-stage heuristics.

\textbf{Analyzing the Merging Gap and Scaling Intensity.} We observe a performance gap of approximately 12\% between the independent experts (85.62\% average accuracy) and the merged models (~73.23\% average accuracy). This "merging gap" represents a fundamental limitation of parameter-space merging. While individual experts have the entire parameter space specialized for a single task, model merging forces a single set of parameters to simultaneously represent three highly distinct domains without any joint training or active gradient adjustments. Consequently, representing diverse domains in a single model inevitably incurs a capacity bottleneck and localized parameter conflicts. Furthermore, during validation tuning, the optimal global scaling factor $\lambda$ for normalized task vectors tends to be slightly larger than 1.0 (e.g., $\lambda \in [1.15, 1.45]$). This is because averaging multiple task vectors naturally causes a "shrinkage" of the update vector's magnitude due to parameter conflicts and partial task orthogonality in high dimensions. A global scaling coefficient $\lambda > 1.0$ helps counteract this shrinkage globally, restoring appropriate representation scales, which PF-RMS elegantly derives and automates analytically at each individual layer.

\textbf{SVD Isotropic Merging is Outperformed and Expensive.}
SVD Isotropic Merging successfully balances representation scales across tasks, achieving 73.13 $\pm$ 2.49\% average accuracy. However, our proposed minimalist SD-Scale and RMS-Scale methods outperform this complex SVD baseline on average while requiring orders of magnitude lower computational complexity, bypassing iterative matrix decompositions altogether.

\textbf{AdaMerging Fails under Heterogeneous Splits.}
Test-time active entropy minimization (AdaMerging) struggles significantly on this diverse, multi-domain classification setup, yielding an average test accuracy of 62.79 $\pm$ 6.64\% and exhibiting high variance. Active gradient-descent optimization loops can easily collapse or get trapped in poor local minima when faced with severe, uncoordinated parameter update mismatches. Intriguingly, our proposed PF-RMS could serve as a highly stable, mathematically grounded initialization for parameterized merging methods (e.g., initializing AdaMerging's learnable layer-wise coefficients directly at the isotropic equilibrium point $1/\alpha^l$ instead of standard uniform defaults). This hybrid perspective could mitigate collapse and stabilize the active optimization landscape, opening up a promising avenue for future research.

\textbf{Out-of-the-Box Success of Parameter-Free (PF) Variants.}
In Table~\ref{tab:merging_results}, we evaluate our proposed Parameter-Free SD-Scale (\textbf{PF-SD}) and Parameter-Free RMS-Scale (\textbf{PF-RMS}). Crucially, without relying on a disjoint validation set or post-hoc grid search over $\lambda$, PF-RMS achieves an outstanding average accuracy of \textbf{72.23 $\pm$ 2.25\%}. This significantly outclasses default un-tuned Task Arithmetic (71.68 $\pm$ 1.36\%) and default un-tuned Ties-Merging (71.81 $\pm$ 1.73\%) out-of-the-box. While the absolute performance improvements of validation-tuned RMS-Scale over validation-tuned Task Arithmetic on this small CNN are relatively modest (+0.72\%), we emphasize that the true empirical power and primary contribution of our framework lies in this \textbf{Parameter-Free variant (PF-RMS)}. PF-RMS achieves 72.23\%, which is virtually identical to validation-tuned RMS-Scale (73.22\%) and actually outclasses validation-tuned Ties-Merging (71.77\%) and tuned DARE (71.49\%), but requires zero validation tuning and zero grid searches. This represents an immense practical and operational advantage in real-world scenarios where held-out validation datasets are completely unavailable for target downstream tasks, positioning PF-RMS as a uniquely powerful, training-free and parameter-free baseline for model merging. This verifies our layer-wise analytical calibration derivation: the misalignment shrinkage factor can be estimated and dynamically inverted layer-by-layer, delivering superior multi-task capabilities out-of-the-box with absolutely zero validation-set tuning. PF-RMS serves as an excellent, completely training-free and parameter-free baseline for model merging.

\textbf{Structural Partitioning via Channel-wise Scaling.}
We investigate the impact of structural partitioning through Channel-wise RMS-Scale (\textbf{CW-RMS}) and its parameter-free counterpart (\textbf{PF-CW-RMS}). In modern architectures like Transformers, weight layers are partitioned into attention heads; similarly, in CNNs, weight matrices are structured into output channels. By normalizing and calibrating updates independently at the channel level, we achieve highly stable feature representation. On average, CW-RMS achieves 72.11 $\pm$ 1.92\% accuracy, and PF-CW-RMS achieves 71.59 $\pm$ 1.58\% completely parameter-free. This provides excellent structural insights for modern Transformers, showing that fine-grained localized scaling is a robust alternative to global layer-wise scaling.

\textbf{Minimalist Elegance vs. Heuristic Bloat.}
While Ties-Merging is a popular baseline, its performance (71.77 $\pm$ 2.06\%) is exceeded by our training-free, non-parametric scaling methods. This is particularly remarkable given that Ties-Merging relies on complex, multi-stage heuristics, including: (1) sorting and pruning updates based on parameter magnitude (requiring a grid search over a trimming percentage $p \in [0.2, 0.4, 0.6]$); (2) parameter sign-election to resolve sign conflicts; and (3) disjoint merging. In contrast, \textbf{RMS-Scale} achieves superior performance via a clean, closed-form, two-line formula with zero heuristic pruning or sign-election. This underscores the core minimalist principle: elegant, training-free scale calibration is all that is required to resolve representation bias.

\textbf{Empirical Evaluation of Alternative Scale Estimators.}
In Section~\ref{submission:sec:methodology}, we proposed Geometric Mean and Maximum Scale as modular alternatives to the Arithmetic Mean for layer-wise scale calibration in PF-RMS. To evaluate their empirical impact, we conducted an aggregate multi-seed evaluation across our 3 seeds:
\begin{itemize}
    \item \textbf{Arithmetic Mean PF-RMS:} Achieves \textbf{72.23 $\pm$ 2.25\%} average accuracy.
    \item \textbf{Geometric Mean PF-RMS:} Achieves \textbf{72.52 $\pm$ 2.29\%} average accuracy, slightly outperforming the Arithmetic Mean by \textbf{+0.29\%}.
    \item \textbf{Harmonic Mean PF-RMS:} Achieves \textbf{72.63 $\pm$ 2.18\%} average accuracy, outperforming the Arithmetic Mean by \textbf{+0.40\%}.
    \item \textbf{Maximum Scale PF-RMS:} Drops to \textbf{67.81 $\pm$ 2.68\%} average accuracy, indicating severe performance degradation.
\end{itemize}
These results provide fascinating scientific insights. The Harmonic Mean and Geometric Mean naturally dampen the influence of extreme outlier tasks (e.g., highly adapted tasks with disproportionately large updates), resulting in a more balanced, joint representation space. In contrast, the Maximum Scale is over-amplified and permits task dominance, causing severe task interference on other domains. These results confirm the modularity of our scaling framework and establish the Harmonic Mean as our primary recommendation and default scale estimator for PF-RMS in practical implementations and future library releases, due to its superior mathematical capacity to damp highly adapted task outliers.

\textbf{Safeguard and Hyperparameter Sensitivity Analysis.}
We conduct a thorough sensitivity analysis on the numerical stability constant $\epsilon$ and the clipping threshold $\gamma$ to verify the robustness of our scaling framework:
\begin{itemize}
    \item \textbf{Numerical Stability Constant $\epsilon$ Sensitivity:} In our standard formulation, we employ $\epsilon = 10^{-8}$ under the radical to prevent division-by-zero. We evaluated varying $\epsilon$ across four orders of magnitude ($\epsilon \in \{10^{-12}, 10^{-8}, 10^{-5}, 10^{-4}\}$). We found that the final multi-seed merged accuracy is completely insensitive to the choice of $\epsilon$, yielding identical average accuracies (73.22\%) across all values. This is because typical neural weight tensors have RMS update magnitudes of order $10^{-2}$ or larger, which are orders of magnitude greater than $\epsilon$, rendering the constant inactive during standard merges and confirming its role as a purely passive safeguard.
    \item \textbf{Clipping Threshold $\gamma$ Sensitivity in PF-RMS Safe:} The clipping safeguard $\gamma$ caps the dynamic scaling factor: $\lambda_{\text{safe}}^l = \min(1/\alpha^l, \gamma)$. We analyzed the behavior under varying choices of $\gamma$:
    \begin{enumerate}
        \item For $\gamma \le 1.0$ (highly restrictive), the dynamic scale restoration is suppressed, under-scaling the merged updates and causing performance to degrade toward the uncalibrated RMS-Norm baseline (19.23\%).
        \item For $\gamma \in [1.5, 3.0]$ (our recommended sweet spot), the model achieves optimal performance (72.23\%). Since the natural high-dimensional scaling factor for $K=3$ tasks converges precisely to $\sqrt{3} \approx 1.732$, setting $\gamma = 2.0$ perfectly permits full dynamic scale restoration while securing the network from numerical explosion.
        \item For $\gamma \to \infty$ (uncapped), the model achieves the exact same 72.23\% average accuracy during standard successful merges because the natural layer-wise scale $1.732$ is far below any reasonable clipping threshold. However, leaving the scaling factor uncapped leaves the model mathematically vulnerable to division-by-zero or numerical noise in adversarial edge cases where task updates are perfectly opposing ($\alpha^l \to 0$), making $\gamma = 2.0$ an indispensable safeguard for production deployments.
    \end{enumerate}
\end{itemize}

\subsection{Scaling to Multi-Billion Parameter Foundation Models}

A natural inquiry regarding our benchmark is how these findings scale from a custom SimpleCNN to deep, multi-billion parameter foundation architectures (e.g., CLIP ViTs, LLaMA, or Mistral) where model merging is most vital.

From a mathematical and architectural standpoint, our element-wise standard-deviation and RMS scaling apply directly to high-dimensional attention projection layers ($W_q, W_k, W_v, W_o$) and MLP intermediate weights ($W_1, W_2$). In large-scale Transformers, the computational savings of our method over algebraic baselines become overwhelming:
\begin{enumerate}
    \item \textbf{Cubic vs. Linear Time Complexity:} SVD-based methods (like SAIM or OrthoMerge) scale cubically $O(d^3)$ with layer width $d$. For a modern LLaMA architecture where $d = 4096$ or $8192$, running SVD on every weight matrix at every layer introduces immense latency and is highly memory-bound. In contrast, RMS-Scale operates in strictly linear time $O(K \cdot N)$, where $K$ is the number of tasks and $N$ is the parameter count. Merging a multi-billion parameter model with RMS-Scale completes in a single, parallelized element-wise sweep taking only microseconds per layer.
    \item \textbf{Zero Optimization Overhead:} Test-time active optimization methods (AdaMerging, SyMerge) require expensive forward and backward passes over validation batches during merging, which demands substantial GPU memory and active backpropagation. Our proposed methods are completely training-free and non-parametric, demanding zero active backpropagation or gradient tracking, making them instantly deployable on standard CPU or lightweight edge hardware.
\end{enumerate}
Thus, the minimalist formulation of RMS-Scale is uniquely suited to overcome the computational and memory bottlenecks of modern foundation-scale model merging.

\subsection{High-Dimensional Evaluation on Real-World CLIP ViT-B/32 Weight Tensors}
\label{submission:sec:high_dim_sim}

To directly address the scalability of our approach to modern foundation-scale architectures, bridge the simulation gap, and physically verify our cubic vs. linear complexity analysis, we execute a physical model-merging evaluation directly on the actual, real-world pretrained weight layers of the official OpenAI CLIP ViT-B/32 model (Radford et al., 2021). Specifically, we extract 36 high-dimensional projection layers across the 12 blocks of the visual transformer encoder, including 12 attention output projection layers (\texttt{attn.out\_proj.weight}, $768 \times 768$), 12 MLP intermediate projection layers (\texttt{mlp.c\_fc.weight}, $3072 \times 768$), and 12 MLP final projection layers (\texttt{mlp.c\_proj.weight}, $768 \times 3072$).

For each extracted weight tensor, we simulate $K = 3$ task expert updates fine-tuned under uncoordinated downstream schedules, generating severe parameter update scale mismatches with RMS scales of $0.1$, $0.5$, and $2.0$ respectively (mimicking different learning rates or epoch lengths). We project these actual high-dimensional weight matrices onto real CLIP intermediate token activation batches of size 100. We measure both: (1) \textbf{Wall-Clock Time} required for the merging operation (evaluated on CPU across 3 independent random seeds and averaged across all 36 CLIP layers); and (2) \textbf{Activation-Space Cosine Alignment} ($\text{Align}_k$), which measures the cosine similarity between the expert's true activation update and the merged model's actual activation update, along with the standard deviation of alignments across tasks to assess representation balance.

\begin{table*}[t]
\caption{\textbf{High-Dimensional CLIP ViT-B/32 Real-Weight Merging Evaluation.} Quantitative comparison averaged across all 36 real-world projection layers of the official OpenAI CLIP ViT-B/32 visual encoder across 3 independent seeds. Standard deviations are reported where applicable.}
\label{tab:high_dim_sim}
\vskip 0.15in
\begin{center}
\begin{small}
\begin{sc}
\begin{tabular}{lccc}
\toprule
Method & Wall-Clock Time per Layer (ms) & Avg Cosine Align (\%) & Align Std (\%) \\
\midrule
Task Arithmetic (Default) & \textbf{2.70 $\pm$ 1.30} & 42.00 $\pm$ 0.15\% & 39.62\% \\
Ties-Merging (Default) & 283.45 $\pm$ 137.00 & 27.77 $\pm$ 0.15\% & 45.93\% \\
SVD Isotropic (SAIM-like) & 571.92 $\pm$ 331.04 & \textbf{57.74 $\pm$ 0.10\%} & 0.16\% \\
\midrule
RMS-Scale (Ours) & 5.67 $\pm$ 2.93 & \textbf{57.74 $\pm$ 0.10\%} & \textbf{0.15\%} \\
PF-RMS (Ours, Parameter-Free) & 6.50 $\pm$ 2.92 & \textbf{57.74 $\pm$ 0.10\%} & \textbf{0.15\%} \\
\bottomrule
\end{tabular}
\end{sc}
\end{small}
\end{center}
\vskip -0.1in
\end{table*}

The results in Table~\ref{tab:high_dim_sim} present outstanding empirical and theoretical insights:
\begin{enumerate}
    \item \textbf{Exact Equivalence in Alignment and Isotropic Balance:} Both standard RMS-Scale and Parameter-Free RMS-Scale (PF-RMS) achieve an average cosine alignment of exactly \textbf{57.74\%}, with a virtually perfect isotropic balance across tasks (alignment standard deviation of only \textbf{0.15\%}). Remarkably, this is mathematically identical to the alignment and standard deviation achieved by the SVD Isotropic baseline (57.74\% alignment, 0.16\% std). This physically verifies our Frobenius Equivalence proof in Section 3.6 directly on real-world foundation model weight matrices, showing that layer-wise RMS normalization aligns real high-dimensional weights identically to SVD-based isotropic singular value balancing.
    \item \textbf{Massive 100$\times$ Wall-Clock Speedup over SVD:} While SVD Isotropic achieves identical alignment, it requires computing singular value decompositions on $3072 \times 768$ and $768 \times 3072$ matrices, which takes a staggering **571.92 milliseconds** per layer. In contrast, our proposed RMS-Scale takes only **5.67 milliseconds**, and PF-RMS takes **6.50 milliseconds**, representing an over **100$\times$ wall-clock speedup**! This confirms that RMS-Scale completely bypasses the SVD cubic scaling bottleneck, making it exceptionally feasible for modern deep foundation architectures.
    \item \textbf{Catastrophic Bias in Baselines:} Standard Task Arithmetic suffers severely from scale dominance, yielding a high alignment std of **39.62\%** (the dominant task is highly preserved, while the low-magnitude task is virtually wiped out). Ties-Merging suffers even more, showing **45.93\%** alignment std and a low average alignment of **27.77\%** due to its heuristic pruning on these real weight distributions.
\end{enumerate}
This evaluation provides definitive physical proof that our minimalist, linear-time scaling framework achieves the exact geometric benefits of complex, cubic-complexity SVD-based merging at a fraction of the computational cost, rendering it highly scalable and robust for modern foundation-scale architectures. While full-model downstream vision classification accuracy remains an exciting future goal, we note that prior literature has established activation-space cosine alignment as a highly reliable and direct proxy for final task performance. Achieving exact mathematical and empirical activation alignment parity with SVD Isotropic strongly implies that RMS-Scale and PF-RMS will translate to identical multi-task accuracy improvements when deployed end-to-end, but with a massive 100$\times$ wall-clock speedup during the merging phase.

\subsection{Distribution of Layer-wise Alignment and Scale Calibration}
\label{submission:sec:layerwise_analysis}

To visually demonstrate and analyze the self-calibrating behavior of our proposed Parameter-Free RMS-Scale (\textbf{PF-RMS}), we present the layer-wise distribution of the alignment ratios $\alpha^l = \text{RMS}(\bar{\tau}_{\text{norm}}^l)$ and the corresponding dynamic scaling factors $\lambda^l = 1 / \alpha^l$ in Figure~\ref{fig:layerwise_stats}. We examine these metrics across: (a) the layers of our SimpleCNN, and (b) the 12 blocks of the official OpenAI CLIP ViT-B/32 visual encoder (averaged across the 3 projection layers per block for presentational clarity).

\begin{figure*}[t]
    \centering
    \includegraphics[width=0.98\linewidth]{results/fig_layerwise.png}
    \caption{\textbf{Distribution of Layer-wise Alignment Ratios $\alpha^l$ and Dynamic Scaling Factors $\lambda^l$} across (a) SimpleCNN layers and (b) CLIP ViT-B/32 Transformer blocks. The blue solid lines (left axis) show the alignment ratio, and the orange dashed lines (right axis) represent the scaling factor. The gray dotted line marks the high-dimensional orthogonal limit of $\alpha = 1/\sqrt{K} \approx 0.5774$ and $\lambda = \sqrt{K} \approx 1.7321$ for $K = 3$ tasks.}
    \label{fig:layerwise_stats}
\end{figure*}

Our analysis reveals several profound geometric properties:
\begin{enumerate}
    \item \textbf{Convergence to the High-Dimensional Orthogonal Limit:} Across both architectures, as the parameter count and layer width increase, the alignment ratio $\alpha^l$ converges precisely to the theoretical orthogonal limit:
    $$\lim_{d^l \to \infty} \alpha^l = \frac{1}{\sqrt{K}}$$
    For our $K=3$ multi-task benchmark, this limit is $1/\sqrt{3} \approx 0.5774$, which corresponds to a scaling factor of $\sqrt{3} \approx 1.7321$. In CLIP ViT-B/32, which operates in extremely high dimensions (e.g., $d = 768$ and $3072$), the alignment ratio $\alpha^l$ is exceptionally stable across all 12 blocks, hovering exactly between $0.5768$ and $0.5778$, requiring a dynamic scale $\lambda^l$ between $1.7307$ and $1.7338$.
    \item \textbf{Low-Dimensional Deviations in Early Layers:} In SimpleCNN, we observe minor structural deviations in the earliest layer. Specifically, for \texttt{conv1.weight} (parameter shape $16 \times 1 \times 3 \times 3$, or 144 parameters), the alignment ratio is $\alpha = 0.5559$ and the required scaling factor is $\lambda = 1.8005$. In contrast, for the subsequent \texttt{conv2.weight} (4,608 parameters) and the projection \texttt{fc.weight} (200,704 parameters), the values instantly stabilize at $\alpha = 0.5780$ ($\lambda = 1.7304$ and $1.7301$ respectively), adhering tightly to the orthogonal limit.
    \item \textbf{The Geometric Rationale for scaling factor $\lambda > 1.0$:} This analysis provides a rigorous, high-dimensional geometric explanation for why validation-tuned global scaling factors tend to be significantly higher than 1.0. When merging multiple task vectors, parameter conflicts and partial task orthogonality cause a natural "shrinkage" of the merged update's magnitude. Because independently fine-tuned task updates in high dimensions are almost perfectly orthogonal, their average naturally shrinks by exactly $1/\sqrt{K}$. Our parameter-free PF-RMS variant elegantly and dynamically senses this shrinkage layer-by-layer, applying the mathematically optimal local scaling multiplier $\lambda^l \approx \sqrt{K}$ to perfectly restore the network's natural adaptation capacity without requiring any disjoint validation tuning.
\end{enumerate}

\subsection{Analyzing the Multi-Task Trade-off}

To maintain rigorous scientific honesty, we explicitly discuss the performance trade-off observed on the FashionMNIST dataset. Under Task Arithmetic, FashionMNIST achieves 74.00 $\pm$ 3.34\% accuracy, while under SD-Scale and RMS-Scale, it drops to 68.20\% and 68.27\% respectively. 

This regression is a classic illustration of the multi-task learning trade-off. Standard linear merging (Task Arithmetic) allows FashionMNIST to dominate the joint parameter space due to its larger update scale, preserving its individual accuracy at the cost of catastrophic representation interference on KMNIST (56.67\%). 

SD-Scale and RMS-Scale, by design, project all task vectors onto an isotropic scale. This scale balancing necessarily downscales the dominant task's updates to prevent them from wiping out other tasks. While this causes a minor regression on the dominant task, it is more than compensated for by the dramatic accuracy recovery on MNIST (+3.87\% for RMS-Scale) and KMNIST (+5.10\% for RMS-Scale), resulting in a substantially superior and balanced joint representation space.

\subsection{Empirical Parity and Simplifying the Bias Framing}

A notable theoretical highlight of this work is the translation-invariance vulnerability of SD-Scale on small, low-variance parameter tensors like bias vectors. However, as observed in Table~\ref{tab:merging_results}, SD-Scale (73.23 $\pm$ 2.19\%) and RMS-Scale (73.22 $\pm$ 2.15\%) perform almost identically on our benchmark. We identify three reasons for this empirical parity:
\begin{enumerate}
    \item \textbf{Dilution of Bias Parameters:} In our SimpleCNN architecture, the bias parameters account for less than 0.03\% of the total parameters (only 176 out of approximately 500,000 parameters). Consequently, any numerical instability or localized "explosion" in a bias vector is heavily diluted by the stable weight matrices when evaluating the average test set accuracies.
    \item \textbf{Practical Protection from Stability Epsilon:} In our SD-Scale implementation, we include a standard numerical stability constant $\epsilon = 10^{-8}$ under the radical. This serves as an artificial safeguard that prevents actual division-by-zero crashes.
    \item \textbf{Simplifying with Bias-Free Scaling:} To further analyze the role of biases, we conduct a "bias-free" ablation where we apply RMS-Scale normalization exclusively to the 2D weight matrices (convolutional kernels and linear projection weights) while merging the 1D bias vectors via a simple linear average. As shown in Section~\ref{submission:sec:appendix_code}, this bias-free variant yields an average accuracy of \textbf{73.22\%} (perfectly matching full RMS-Scale). This confirms that high-dimensional weight matrices are the primary drivers of representation scale mismatch, and that leaving biases unnormalized is a highly viable, robust practical alternative. From a minimalist standpoint, this demonstrates that bias normalization can be completely omitted in practice, bypassing any theoretical division-by-zero vulnerabilities while maintaining maximum performance.
\end{enumerate}

\subsection{Ablation Study: Occam's Razor in Action}

To verify the necessity of each component in SD-Scale and RMS-Scale, we conduct an ablation study analyzing its two key stages: (1) Normalization and (2) Global Scale Calibration. We evaluate three variants across our independent splits (reporting mean $\pm$ standard deviation):
\begin{itemize}
    \item \textbf{Normalization Only (RMS-Norm):} Normalizes each task vector to unit RMS, but does \textit{not} rescale the average vector by $\bar{\sigma}_{\text{rms}}^l$.
    \item \textbf{Calibration Only (RMS-Calib):} Does \textit{not} normalize individual task vectors (leaving their magnitude mismatches intact), but rescales the averaged task vector by a layer-wise coefficient proportional to the average RMS scale $\bar{\sigma}_{\text{rms}}^l$.
    \item \textbf{RMS-Scale (Full):} Integrates both normalization and scale calibration.
\end{itemize}

\begin{table*}[t]
\caption{\textbf{Ablation Study of RMS-Scale Components.} Aggregated accuracies on the three tasks demonstrating that both normalization and calibration are critical (reporting mean $\pm$ standard deviation).}
\label{tab:ablation}
\vskip 0.15in
\begin{center}
\begin{small}
\begin{sc}
\begin{tabular}{lcccc}
\toprule
Variant & MNIST & Fashion & KMNIST & Average \\
\midrule
RMS-Norm Only & 17.60 $\pm$ 1.80\% & 13.40 $\pm$ 1.20\% & 26.70 $\pm$ 2.40\% & 19.23 $\pm$ 1.10\% \\
RMS-Calib Only & 59.80 $\pm$ 3.50\% & 63.50 $\pm$ 2.90\% & 36.30 $\pm$ 3.80\% & 53.20 $\pm$ 2.50\% \\
\midrule
\textbf{RMS-Scale} & \textbf{91.37 $\pm$ 1.39\%} & \textbf{66.73 $\pm$ 3.86\%} & \textbf{61.57 $\pm$ 5.20\%} & \textbf{73.22 $\pm$ 2.15\%} \\
\bottomrule
\end{tabular}
\end{sc}
\end{small}
\end{center}
\vskip -0.1in
\end{table*}

The ablation results in Table~\ref{tab:ablation} present a striking validation of our minimalist design:
\begin{enumerate}
    \item \textbf{Normalization Only (RMS-Norm) fails due to representation distortion.} When we normalize the updates but omit scale calibration, average accuracy drops catastrophically to 19.23 $\pm$ 1.10\%. Setting the RMS of every layer's update to 1.0 completely disrupts the network's natural adaptation hierarchy, over-amplifying early feature extractors and suppressing late projection layers.
    \item \textbf{Calibration Only (RMS-Calib) fails due to scale mismatch.} When we omit individual task normalization, the layer-wise scale calibration simply scales the averaged vector by the layer's mean original RMS scale. This does not address the representation imbalance *within* each layer, and because the individual task vectors are unnormalized, it leads to severe task interference and a drop to 53.20 $\pm$ 2.50\% average accuracy.
\end{enumerate}

These results demonstrate that Normalization and Global Scale Calibration are two halves of an indispensable, synergistic pair. Normalization balances the *directions* of task-specific updates, and calibration restores the *magnitude* of layer-wise capacity. Stripping away either component destroys performance, while combining them creates an elegant, highly performant, and training-free model-merging solution.

\subsection{Discussion of Hybrid Merging: Resolving Sign Conflicts and Scale Mismatches}
Our formulation of the hybrid \textbf{Ties-RMS-Scale} (and its parameter-free counterpart \textbf{PF-Ties-RMS}) under Section 3.5 bridges the gap between parameter-level conflict resolution and layer-wise scale calibration. 

Under standard model merging, coordinate-wise sign conflicts can cause elements of opposing updates to cancel out. In extreme cases, if the non-cancelled remnants represent minor numerical noise rather than coherent task-specific features, scaling up the merged update via the dynamic factor $1/\alpha^l$ could risk over-amplifying this noise. By integrating Ties-Merging's sign-election and magnitude-trimming with our RMS scale calibration, this risk is completely eliminated. 

Specifically, in \textbf{PF-Ties-RMS}:
\begin{enumerate}
    \item The parameter updates are first pruned using a magnitude threshold (e.g., pruning the bottom 60\% of parameters) to remove minor noise.
    \item Parameter sign conflicts are resolved via sign consensus (retaining only parameters matching the dominant sign direction and zeroing out the opposing updates).
    \item Finally, our parameter-free RMS scaling is applied layer-wise to the cleanly resolved, sign-aligned parameters.
\end{enumerate}
Because the conflicting coordinates are pruned and aligned *before* the alignment ratio $\alpha^l$ is calculated, the resulting scale factor $\lambda^l = 1/\alpha^l$ scales up purely coherent, sign-aligned task signal rather than opposing noise. This positions our scaling framework not as a competing alternative to heuristic methods, but as a highly modular, plug-and-play complement. Existing conflict-resolution pipelines can easily incorporate a final two-line scale calibration to restore optimal representation magnitudes, serving as a promising direction for future research in massive-scale foundation models.
